% !Mode:: "TeX:UTF-8"

\chapter{数据采集与特征工程}

\section{采集流程}

项目通过 Python 脚本完成从 Prompt 生成到模型回复采集的流程。核心文件包括：

\begin{itemize}
    \item \texttt{scripts/01\_generate\_prompts.py}：生成 Prompt 矩阵；
    \item \texttt{scripts/02\_collect\_responses.py}：调用 mock 或真实模型 API；
    \item \texttt{scripts/03\_build\_features.py}：从回复文本中抽取行为特征；
    \item \texttt{scripts/04\_analyze.py}：执行统计分析、聚类和关联规则挖掘。
\end{itemize}

采集脚本支持断点续采，能够跳过已经存在的 $(model, repeat, prompt\_id)$ 组合，避免接口调用重复。真实模型采集时，应记录模型名称、采集时间、temperature、最大输出长度和失败重试情况，保证实验可追溯。

\section{数据文件说明}

项目中的主要数据文件如表 \ref{tab:data-files} 所示。

\begin{table}[htbp]
    \centering
    \bicaption{项目主要数据文件}{Main data files of the project}
    \begin{tabular}{ll}
        \toprule
        文件路径 & 内容说明 \\
        \midrule
        \texttt{data/prompts.csv} & Prompt 实验矩阵 \\
        \texttt{data/raw/model\_responses.csv} & 模型原始回复 \\
        \texttt{data/processed/behavior\_features.csv} & 行为特征表 \\
        \texttt{outputs/tables/*.csv} & 统计表、聚类结果和关联规则 \\
        \texttt{outputs/figures/*.png} & 热力图、雷达图和聚类图 \\
        \bottomrule
    \end{tabular}
    \label{tab:data-files}
\end{table}

\section{行为特征抽取}

模型原始回复通常是多段自然语言文本，无法直接进行统计分析。因此，需要将回复转换为结构化行为特征。本文关注的主要特征如表 \ref{tab:features} 所示。

\begin{table}[htbp]
    \centering
    \bicaption{模型回复行为特征}{Behavioral features extracted from model responses}
    \begin{tabular}{ll}
        \toprule
        字段 & 含义 \\
        \midrule
        ai\_judgment & 模型最终判断，包含能通过、不能通过、不确定 \\
        is\_correct & 模型判断是否与参考标签一致 \\
        reasoning\_steps & 回复中的推理步骤数量 \\
        has\_formula & 是否出现公式、比较符号或变量表达 \\
        uses\_coordinate & 是否使用坐标系、包络线等建模语言 \\
        error\_category & 错误回复对应的失效模式类别 \\
        \bottomrule
    \end{tabular}
    \label{tab:features}
\end{table}

\section{失效模式定义}

本项目将错误回复初步归纳为以下几类：概念混淆、计算崩溃、直觉过拟合、噪声牵引、单位混乱和信息不足。后续正式报告可结合典型样例，说明每类错误在文本中的具体表现。
